% Created 2018-06-01 Fri 13:38
% Intended LaTeX compiler: pdflatex
\documentclass[11pt]{article}
\usepackage[utf8]{inputenc}
\usepackage[T1]{fontenc}
\usepackage{graphicx}
\usepackage{grffile}
\usepackage{longtable}
\usepackage{wrapfig}
\usepackage{rotating}
\usepackage[normalem]{ulem}
\usepackage{amsmath}
\usepackage{textcomp}
\usepackage{amssymb}
\usepackage{capt-of}
\usepackage{hyperref}
\usepackage[margin=0.80in]{geometry}
\usepackage[backend=biber, style=ieee, url=false]{biblatex}
\usepackage{float}
\usepackage[super,negative]{nth}
\usepackage[capitalise]{cleveref}
\usepackage{pst-node,transparent,ragged2e}
\addbibresource{/home/nuk3/.spacemacs.d/org-files/bibliography.bib}
\DeclareFieldFormat[inproceedings]{citetitle}{\textit{#1}}
\DeclareFieldFormat[inproceedings]{title}{\textit{#1}}
\DeclareFieldFormat[misc]{citetitle}{#1}
\DeclareFieldFormat[misc]{title}{#1}
\renewcommand*{\bibpagespunct}{%
\ifentrytype{inproceedings}
{\addspace}
{\addcomma\space}}
\AtEveryCitekey{\ifuseauthor{}{\clearname{author}}}
\AtEveryBibitem{\ifuseauthor{}{\clearname{author}}}
\SpecialCoor
\author{Nyameko Lisa}
\date{18 May 2018}
\title{IPSP066 - Internet Aspects of Copyright and Trade Marks\\\medskip
\large Assignment 3: 827826}
\hypersetup{
 pdfauthor={Nyameko Lisa},
 pdftitle={IPSP066 - Internet Aspects of Copyright and Trade Marks},
 pdfkeywords={},
 pdfsubject={},
 pdfcreator={Emacs 27.0.50 (Org mode 9.1.2)},
 pdflang={English}}
\begin{document}

\maketitle
\addvspace{110pt}
\centering{
\pnode(0.5\textwidth,-0.5\textheight){thisCenter}
\rput(thisCenter){%\transparent{0.25}
\includegraphics[width=2.7in]{/home/nuk3/course/llb/wipo-unisa/UNISACoatofArms.eps}}}

\justifying
\addvspace{110pt}
\section*{Declaration}
\label{sec:org4324d9c}
I know that plagiarism is to use someone else’s work and pass it off as my own.\\
I know that plagiarism is wrong.\\
I confirm that this assignment is my own work.\\
I have acknowledged in the bibliography accompanying the assignment all the sources that I have used.\\
I have not directly copied without acknowledgement anything from the Internet or from any other source.\\
I have indicated every quotation and citation in a footnote or bracket linked to that quotation.\\
I have not allowed anyone else to copy my work and to pass it off as their own work.\\
I understand that if any unacknowledged copying whatsoever appears in my assignment I will receive zero per cent for the assignment.\\
I am aware of the UNISA policy on plagiarism and understand that disciplinary proceedings can be instituted against me by UNISA if I contravene this policy.\\
I indicate my understanding and acceptance of this declaration by
entering my name here under:
\begin{itemize}
\item Name: \textbf{Nyameko Lisa} (Student Number: \textbf{7874-909-3})
\end{itemize}

\subsection*{NOTE}
\label{sec:org02e5544}
Please note that footnotes will be denoted as \footnote{This is a footnote.} and will
appear at the bottom of the page.\\
References will be denoted by \cite{rsa93_tm_act} and will appear at the end of the document.
\newpage

\section{Give a list of fifteen (15) questions that you will ask Sibusiso in order to advise him in this matter.}
\label{sec:org05c4ce7}
First Sibusiso will be advised on the aspects of his domain name as they pertain
to trade mark law. Then Sibusiso will be advised in terms of the content of his
web site as it pertains to the law of copyright. Finally, a course of action for
Sibusiso to pursue will be presented based on a summative evaluation
of his circumstances.
\subsection{What is the `nature' of the Sound Recordings Sibusiso wishes to sell on his website?}
\label{sec:org3e0902b}
As per the definitions of [section 2(1)]\cite{rsa93_tm_act} and given that the
Republic is a Paris Convention Member State, [article 15(1)]\cite{wto17_trips},
[article 1(2)]\cite{wipo83_paris_conve_protect_ip}, it follows that
\emph{`UNITED.gTLD'} and \emph{`UNITED.ccTLD} will be understood to be \textbf{marks}, in that
they are names or signs capable of being represented graphically. Moreover
\emph{`UNITED.COM'} constitutes a \textbf{well-known registered service} or \textbf{trade mark} in
the United States of America in relation to distinguishing their trade or use
from other goods or services connected in the course of trade regarding
\emph{`airline and courier services.'}, as per the provisions of [Class 39 -
Transport; packaging and storage of goods, travel
arrangements]\cite{wipo57_ncl}. As per those very same provisions, Sibusiso would be advised to pursue
trade mark registration of the `united' mark under:
\begin{itemize}
\item Class 35: advertising; business management/administration - \emph{``in particular the bringing together, for the benefit of others, a variety of goods enabling customers to conveniently view and purchase those goods; such services may be provided by\ldots{} means of electronic media, for example, through web sites\ldots{}''}, and
\item Class 41: education; providing of training; entertainment; sporting and
cultural activities.
\end{itemize}

\subsection{Has Sibusiso already approached an ICANN-Accredited reseller to query the availability of UNITED.COM?}
\label{sec:orge6fee01}
In order to reserve such a domain name in a general Top-Level Domain or \textbf{gTLD},
the \textbf{registrant} Sibusiso must register it with an ICANN-accredited
registrar\footnote{Or alternatively through a registrar's resellers.}
\cite{icann13_rules}. The registrar will then check if the UNITED.gTLD domain name
is available, and then either create or lookup a WHOIS (domain name registration
record) with the registrant's information \cite{icann99_policy}. The registrant would be advised
against the pursuit of gTLD of the `united' domain name, as he is not likely to
be successful. In particular, while the `united' trade mark may be territorial by
its very nature, the specific `united.com' domain name cannot be owned by
multiple companies, irrespective of the uniqueness of their products and/or
markets.

\subsection{Has Sibusiso considered ccTLD registration or an alternative domain name?}
\label{sec:orgb483336}
The registrant would be advised to instead pursue registration of the country
code Top-Level Domain or \textbf{ccTLD} domain `united.co.za', through one of the ccTLD
operators [section 59]\cite{rsa02_elect_comm_trans_act}, particularly if the
students he intends to sell sound recordings to are based within the Republic.
This will be assigned on a first-come first-serve basis. Alternatively, he could
pursue registration of a more descriptive domain name, such as
\uline{`studentsoundz.co.za'}, for example.

\subsection{Does Sibusiso understand that he could be liable for trade mark infringement?}
\label{sec:org0924391}
As per the provisions of [sections 10(6) and 35]\cite{rsa93_tm_act} \emph{(for
litigation within the Republic)}, [Article 16 subsections (1), (2) and
(3)]\cite{wto17_trips}, [Article
6\textsuperscript{\textit{bis}}(1)]\cite{wipo83_paris_conve_protect_ip},
[Article 16]\cite{wipo70_tlt} and [subsection 1125(c)]\cite{usa46_title15_ch22}
\emph{(for actions within the United States of America)}, the proprietors of the
second level domain `united', as is the case with the airline company, may argue
that it is indeed a \textbf{well-known} mark and that registration of Sibusiso's domain
name in either a gTLD or ccTLD, constitutes trade mark infringement, through the
\textbf{dilution} of a well-known mark.

\subsection{Has Sibusiso already began trading under the `united' and has he in any way exploited the `good will' or `brand' of the UNITED.COM airline company?}
\label{sec:orgc154e15}
Dilution protection is not subject to the same limitations of `traditional'
trade mark infringement,\footnote{Restricted to goods and services which are the same or similar to those for which the mark is registered.\label{org03cdd66}} in that it extends to use of the offending mark to
\emph{any} goods and services. The proprietors of the UNITED.com Airlines trade
mark and domain name, need demonstrate that:
\begin{itemize}
\item Their mark is well known within the Republic, i.e. it would be sufficient for
them to show that their trade mark has acquired a reputation amongst a
substantial number of members of the public,
\citetitle{corbett97_mcd_v_joburgers} \cite{corbett97_mcd_v_joburgers},
\item Sibusiso has used the mark in the course of trade, where the courts of
the Republic view this in the `traditional' trade mark infringement\textsuperscript{\ref{org03cdd66}}
sense, \citetitle{harms01_cowbell_v_ics} \cite{harms01_cowbell_v_ics}, lastly
\item Use of the infringing mark is likely to take unfair advantage of, or be
detrimental to the distinctive character of the UNITED Airlines trade mark,
amounting to unfair competition [articles 2(1) and 3(2)(a)(i)]\cite{wipo96_model_provi_unfair_comp}.
\end{itemize}

\subsection{Does Sibusiso know that there are two additional procedures that could be instituted to resolve the domain name dispute?}
\label{sec:orgbc46e53}
ICANN provides a \textbf{Uniform Domain Name Dispute Resolution Policy} which outlines
provisions [section 4(a)(i)]\cite{icann99_policy}, within which mandatory
administrative proceedings may be instituted, on the basis that as a \textbf{third
party complainant} UNITED Airlines asserts that Sibusiso's domain is identical
or confusingly similar to their trade or service mark over which they have
rights to.\\

As per the provisions of [Chapter 2, Paragraphs 110 -
111]\cite{wipo99_management_dns_ip_report}, Sibusiso may institute Submission to
Jurisdiction and to \textbf{Alternative Dispute Resolution} Procedures. The ADR procedures to
follow are described and outlined in Chapter 3: Resolving Conflicts in a
Multijurisdictional World with a Global Medium: A Uniform Dispute-Resolution
Policy.

\subsection{Does Sibusiso know that copyright subsists uniquely in both the samples as well as the original sound recordings he wishes to host for sale on his web site?}
\label{sec:org70d2f39}
As a Convention Member the provisions of [section 5]\cite{rsa78_copyrightact},
[section 104(c,d)]\cite{usa76_title17_us_copyright_act}, [Article
9(1)]\cite{wto17_trips}, [Article 1(2)]\cite{eurlex00_elec_commerce}, [Article
5(2)]\cite{wipo86_berne} and [Article 1(4)]\cite{wipo96_copyright_treaty} which
dictate that the statues of the Republic shall govern international copyright
protection as it pertains to the Internet. Furthermore as per the provisions of
[sections 1 and 2(1)(b,e)]\cite{rsa78_copyrightact}, [sections 101 and
102(a)(2,7)]\cite{usa76_title17_us_copyright_act} and [Article
2(1)]\cite{wipo86_berne} the sound recordings shall be considered as \textbf{protected
artistic works} or \textbf{sound recordings}, where the samples of those works
constitute \textbf{adaptations} or \textbf{derivative works}.

\subsection{Has Sibusiso identified and indicated his intent to sell the sound recordings/samples to the respective owners, i.e. to the person(s) (i) who first made or created the works, (ii) by whom the arrangements for the production of the sound recording were made, or (iii) the first broadcasters of the works?}
\label{sec:orgf7ea0e9}
As per the provisions of [sections 1, 3, 4(1)(a)]\cite{rsa78_copyrightact},
[section 101]\cite{usa76_title17_us_copyright_act}, [Article 14]\cite{wto17_trips},
[Articles 2 and 3]\cite{wipo96_wppt}, [Article 3 and 9(1)]\cite{wipo86_berne}, the
\textbf{authors}, will in this instance be considered as the intellectual property
rights holders of the sound recordings and corresponding samples, and they are
understood to be the performers, producers, licensees and broadcasting
organisations of the sound recordings.

\subsection{Without explicit permission, does Sibusiso understand that he may be in violation of a number of the intellectual property rights that vest with the holders discussed above?}
\label{sec:org887eb84}
\begin{itemize}
\item Exclusive right of reproduction [section 9(a)]\cite{rsa78_copyrightact},
[section 106(1)]\cite{usa76_title17_us_copyright_act}, [Article 9(1) and
(2)]\cite{wipo86_berne}, [Article 7 and 11]\cite{wipo96_wppt} and [Article
2(a-e)]\cite{eurlex01_harmonisation},
\item Exclusive right of offering by way of trade, directly or indirectly, a
reproduction of the sound recording [section 9(b)]\cite{rsa78_copyrightact}, [section 106(3)]\cite{usa76_title17_us_copyright_act},
\item Exclusive right of publication, distribution or broadcast to the public, or
causing the sound recording to be transmitted in a diffusion service [section
9(c-d)]\cite{rsa78_copyrightact}, [Article
11\textsuperscript{bis}(1)(i)-(ii)]\cite{wipo86_berne}, [Article
6(1)]\cite{wipo96_copyright_treaty}, [Article 8(1) and 12(1)]\cite{wipo96_wppt}
and [Article 4]\cite{eurlex01_harmonisation},
\item Exclusive right of communication to the public of a performance [section
9(e)]\cite{rsa78_copyrightact}, [section
106(6)]\cite{usa76_title17_us_copyright_act}, [Article
11(1)(i)-(ii)]\cite{wipo86_berne}, [Article 8]\cite{wipo96_copyright_treaty},
[Article 13-15]\cite{wipo96_wppt} and [Article 3]\cite{eurlex01_harmonisation},
and
\item Exclusive right\footnote{Sampling constitutes an adaptation.} of adaptation,
arrangement and other alteration [section 6(f-g)]\cite{rsa78_copyrightact},
[Article 12]\cite{wipo86_berne} and [Article 16]\cite{wipo96_wppt}.
\end{itemize}

\subsection{Does Sibusiso know how old the sound recordings and samples are, or can he demonstrate that his use amounts to fair dealing, and does he appreciate than it may be necessary for him to become a licensee of the various copyright holders?}
\label{sec:orgb806403}
As per the provisions of [section 3(2)(a,c,d,e,f)]\cite{rsa78_copyrightact},
[Article 7]\cite{wipo86_berne} and [Article 17]\cite{wipo96_wppt} the duration for which copyright subsists in the
sound recordings and their corresponding samples, is \textbf{fifty years} from the end
of the year of first performance, broadcast or publication. Should none of these
apply then the term of protection is the life of the author and fifty years from
the year in which the author dies. Or should infringement action be heard in the
United States, \textbf{95} years after first publication, \textbf{120} years after creation or
\textbf{70} years after the presumed death of the author, whichever expires first as per
[section 302(e)]\cite{usa76_title17_us_copyright_act}.\\

As per the provisions of [section 12]\cite{rsa78_copyrightact}, [section
107]\cite{usa76_title17_us_copyright_act} and [Articles 10 and
10\textsuperscript{bis}]\cite{wipo86_berne}, if he can demonstrate that
his use amounts to \emph{`fair dealing'}, e.g. as per [subsection (c)(ii)], do the
sound recordings constitute the broadcasting for the purpose of reporting
current events, for example?\\

Unless Sibusiso can show that his intended use satisfies either of the above
conditions he will require the authors' permission by way of a license, [section
22]\cite{rsa78_copyrightact}, [section 115]\cite{usa76_title17_us_copyright_act} and
[Article 13(1)]\cite{wipo86_berne}.

\subsection{Has Sibusiso considered the implications of his potential intellectual property rights infringement on his intermediary service provider?}
\label{sec:orgd1e7dd0}
With respect to Sibusiso's intellectual property rights infringement, the
intermediary service provider HostNet shall \uline{\textbf{not}} be legally liable, for
activities conducted during the course of its trade:
\begin{itemize}
\item Information transmission or acting as a \textbf{`mere conduit'} [Article 12]\cite{eurlex00_elec_commerce},
\item Temporary information storage \textbf{`caching'} [Article 13]\cite{eurlex00_elec_commerce}, and
\item Hosting [Article 14]\cite{eurlex00_elec_commerce}.
\end{itemize}

Moreover there exist provisions against statutory requirements obliging HostNet
to monitor Sibusiso's activities on their respective servers [Article 15]\cite{eurlex00_elec_commerce}.

\subsection{Has Sibusiso arranged a Royalty payment agreement or can he demonstrate that he is a `manufacturer' of the sound recordings which are hosted on American Servers?}
\label{sec:org00f106d}
Linking is a means of exploiting third-party content on the Internet, and thus
enhancing a web-page with coloured or underlined descriptive words \textbf{Hyperlinked}
to Uniform Resource Locator (URL), linking to the audio samples. The mere
creation of the links does not in and of themselves constitute copyright
infringement.\\

However, depending on the licensing/contractual conditions of the links to the
samples, Sibusiso may find himself liable to the proprietors and content
generators of the sites hosted on servers in the United States of America on
which the samples appear, \textbf{in addition to} the original authors of the sound
recordings for contributory infringement, should the sound recordings be hosted
on those American servers illegally.\\

As per the provisions of [section 9A]\cite{rsa78_copyrightact} and [section
111]\cite{usa76_title17_us_copyright_act}, in the absence of licensing and/or
contractual agreements to the contrary, Sibusiso will be required to pay
appropriate royalties to the owners of the relevant copyright.\\

Lastly should Sibusiso be capable of demonstrating that he satisfies the
provisions of [section 14(1)(a-d)]\cite{rsa78_copyrightact} and [section
115]\cite{usa76_title17_us_copyright_act} then he meet the requirements of a
``manufacturer'', for which copyright in a musical work shall not be infringed
by such a person.

\subsection{Does Sibusiso appreciate that an infringement action against him may be instituted within the Republic, or within the United States, or even in both jurisdictions?}
\label{sec:org8b30c78}
Infringement action can be instituted against within the Republic or within the
United States of America. In particular however, given that the defendant
Sibusiso is domiciled or resident of, and his business is registered within the
Republic, the High Courts of the Republic have jurisdiction.\\

However should Sibusiso, acting as the defendant in some such matter, be a
foreign peregrinus (neither domiciled nor resident) of all South African High
and Magistrate Courts, then the plaintiff's area of residence, or where the
conclusion and/or breach of contract occurred, i.e. the United States of
America, will hold jurisdiction. Regardless however, the parties cannot
themselves submit to or choose a specific magistrate's courts in their
agreement.

\subsection{Does Sibusiso understand that he may be liable for vicarious liability or contributory infringement?}
\label{sec:org7986eac}
If it can be shown that the users of Sibusiso website, themselves were actively
engaged in the unauthorised reproduction and distribution of the copyrighted
works, both Sibusiso and the users of his website, may find themselves liable
for contributory and vicarious copyright infringement [section
23(2)(b,c,d)]\cite{rsa78_copyrightact} and [sections 110, 501 and
1201]\cite{usa76_title17_us_copyright_act}, in matters to be potentially heard in
jurisdictions across the globe.\\

Within the context of American case law, in the matter heard before the United
States Court of Appeals in the Ninth Circuit, between
\citetitle{scroeder01_anm_records_v_napster}
\cite{scroeder01_anm_records_v_napster}, the onus was on the industry to
demonstrate that Napster users themselves were actively engaged in the
unauthorised reproduction and distribution of copyrighted works, for successful
action of contributory and vicarious copyright against Napster, whose
counterclaim against their users' direct infringement, and in turn themselves,
was the `fair use' defence. The Court of Appeal concurred with the district
court's finding, in that:
\begin{quote}
\textit{``...more than eighty percent of the files available on Napster
may have been copyright protected and owned or administered by plaintiffs,
thereby violating the plaintiff's exclusive rights of reproduction and
distribution under §106, U.S. Copyright Act.}
\end{quote}

\begin{quote}
\textit{The district court determined that Napster had harmed the music industry's
market in two ways: it reduces audio CD sales among college students and it
raises barriers to plaintiffs' entry into the market for the digital downloading
of music. To which the Court of Appeal again concurred, adding the that lack of
harm to an established market cannot deprive the copyright holder of the right
to develop alternative markets for their works.''}
\end{quote}

\subsection{Does Sibusiso understand all the statutory laws that may apply before he proceeds with registration of his domain and commencement of his business trading?}
\label{sec:org7ca03aa}
Issues of contract or trade practice law or both may arise in addition to issues
of trade mark and copyright law. Depending on the "click-wrap" or implied license, a number of
subjective implications may arise pertaining to the whether Sibusiso was in fact
eligible or not to hyperlink the multimedia material for commercial purposes on
his own website. Moreover it is not entirely clear which jurisdiction these
matters may be heard.\\

Sibusiso would be strongly advised to:
\begin{itemize}
\item pursue domain name and trade mark registration of a \textbf{`unique'} ccTLD such as
\emph{studentsounds.co.za} or \emph{edurecordings.co.za}, for example,
\item obtain licensing / contractual authorisations from the \textbf{authors} and copyright holders to host
the sound recordings and their corresponding samples,
\item obtain licensing / contractual authorisation from the American web hosts to use
their samples on his own web site,
\item establish licensing / contractual agreements with his web host to absolve him
of any liability that may arise from his users' (the students) misuse of his platform.
\end{itemize}
\newpage
\printbibliography
\end{document}